\section{Cylinder uniqueness closes the whole sequence}
\label{sec:identification}

Every weak subsequential limit $\nu$ has the target marginal $(\Psi_K)_{\#}G_K$ for each finite mode list $K$. Indeed, the coordinate maps are continuous, so weak convergence on the carrier implies weak convergence of their pushforwards. The full sequence of pushforward laws already converges to the target by Section~\ref{sec:marginals}; uniqueness of limits in the finite-dimensional probability-measure space identifies the pushforward of $\nu$ with that same target.

It remains to show that these finite marginals determine the carrier law. The coordinate maps generate the cylinder pi-system. Every such cylinder is Borel. Separability of the weighted $\ell^2$ carrier makes it Polish, and the formal development proves that the sigma-algebra generated by finite-coordinate cylinders equals the carrier Borel sigma-algebra. The finite-measure uniqueness theorem therefore implies that any two Borel probability measures with the same finite weighted Fourier marginals coincide.

Choose one Prokhorov limit and call it $\mu_{\mathrm{free}}^{L,M}$. The preceding argument shows that every extracted limit has the same finite marginals as this measure, hence equals it. To see that the original sequence converges, suppose it did not. Then some neighborhood of $\mu_{\mathrm{free}}^{L,M}$ would be missed along a subsequence. Tightness would give a further weakly convergent subsequence; its limit must equal $\mu_{\mathrm{free}}^{L,M}$, contradicting that it stays outside the neighborhood. Thus the entire sequence converges.

This step is stronger than identifying one convergent subsequence. It uses the fact that every subsequence admits a further limit and that the Borel cylinder marginals leave no freedom in that limit. The resulting $\mu_{\mathrm{free}}^{L,M}$ is an actual Borel probability measure, and its construction and uniqueness are both part of Theorem~\ref{thm:main}.
